% Part:sets-functions-relations
% Chapter: sets
% Section: reduction

\documentclass[../../../include/open-logic-section]{subfiles}

\begin{document}

\olfileid{sfr}{siz}{red}

\olsection{Reductie: een probleem tot een ander herleiden}

\begin{editorial}
  In deze paragraaf bewijzen we door reductie---door het ene probleem
  tot het andere te herleiden---dat een verzameling niet kan worden
  opgesomd. De uitkomst sluit aan bij de resultaten in
  \olref[nen]{sec}. Een andere, iets kortere versie met dezelfde
  uitkomst als in \olref[nen-alt]{sec} staat in \olref[red-alt]{sec}.
\end{editorial}

Met een diagonaalargument hebben we laten zien dat $\Pow{\PosInt}$
!!{nonenumerable} is. We hadden al een bewijs dat $\Bin^\omega$, de
verzameling van alle oneindige rijen met $0$'en en $1$'en,
!!{nonenumerable} is. We kunnen ook op een andere manier bewijzen dat
$\Pow{\PosInt}$ !!{nonenumerable} is: toon aan dat \emph{als
$\Pow{\PosInt}$ !!{enumerable} is, dan $\Bin^\omega$ ook
!!{enumerable} is}. We weten dat $\Bin^\omega$ niet !!{enumerable} is,
dus $\Pow{\PosInt}$ kan het ook niet zijn. Dit noemen we het ene
probleem tot het andere \emph{reduceren}. Hier herleiden we de vraag
hoe $\Bin^\omega$ kan worden opgesomd tot de vraag hoe
$\Pow{\PosInt}$ kan worden opgesomd. Een oplossing voor die tweede
vraag---een opsomming van $\Pow{\PosInt}$---zou meteen een oplossing
opleveren voor de eerste vraag---een opsomming van $\Bin^\omega$.

Hoe herleiden we het probleem om een verzameling~$B$ op te sommen tot
het probleem om een verzameling~$A$ op te sommen? We moeten van een
opsomming van~$A$ een opsomming van~$B$ kunnen maken. De makkelijkste
manier is een surjectieve functie $f\colon A \to B$ te definiëren,
dus een functie die elke mogelijke uitkomst bereikt. Als de lijst $x_1$,
$x_2$, \dots{} de verzameling~$A$ opsomt, dan zouden $f(x_1)$, $f(x_2)$,
\dots{} de verzameling~$B$ opsommen. In ons geval zoeken we een
surjectieve functie $f\colon \Pow{\PosInt} \to \Bin^\omega$.

\begin{prob}
Stel dat er een injectieve functie $g\colon B \to A$ is, dus een
functie die verschillende invoeren nooit dezelfde uitkomst geeft, en
dat $B$~!!{nonenumerable} is. Toon aan dat~$A$ dan ook
!!{nonenumerable} is. Laat daarvoor zien hoe je met~$g$ een opsomming
van~$A$ kunt omzetten in een opsomming van~$B$.
\end{prob}

\begin{proof}[Bewijs van {\olref[nen]{thm:nonenum-pownat}} via herleiding]
Stel dat $\Pow{\PosInt}$ toch !!{enumerable} zou zijn. Dan zouden we
haar kunnen opsommen: $Z_{1}$, $Z_{2}$, $Z_{3}$, \dots

Definieer de functie $f \colon \Pow{\PosInt} \to \Bin^\omega$ als
volgt. We stellen $f(Z)$ gelijk aan de rij $s_{k}$ waarbij $s_{k}(n) =
1$ precies als $n \in Z$; in alle andere gevallen is $s_k(n) = 0$.
Dit levert echt een functie op. Als $Z \subseteq \PosInt$, is elke $n
\in \PosInt$ namelijk ofwel !!a{element} van $Z$, ofwel niet. Zo wordt
de verzameling $2\PosInt = \{2, 4, 6, \dots\}$ van positieve even
getallen gekoppeld aan de rij $010101\dots$. De lege verzameling wordt
gekoppeld aan $0000\dots$ en de verzameling $\PosInt$ zelf aan
$1111\dots$.

De functie is ook !!{surjective}: ze bereikt elke mogelijke uitkomst.
Elke rij van $0$'en en $1$'en hoort bij een verzameling positieve gehele
getallen. Die verzameling bevat precies de posities
waarop de rij een~$1$ heeft. Om dit precies op te schrijven, neem
$s \in \Bin^\omega$ en definieer $Z \subseteq \PosInt$ door:
\[
Z = \Setabs{n \in \PosInt}{s(n) = 1}
\]
Dan geldt $f(Z) = s$; dat volgt uit de definitie van~$f$.

Bekijk nu de lijst
\[
f(Z_1), f(Z_2), f(Z_3), \dots
\]
Omdat $f$ !!{surjective} is, komt elk element van $\Bin^\omega$ voor
als waarde van~$f$ bij een of ander argument. Elk element staat dus op
de lijst. Deze lijst somt daarom heel~$\Bin^\omega$ op.

Als $\Pow{\PosInt}$ !!{enumerable} was, zou $\Bin^\omega$ dus ook
!!{enumerable} zijn. Maar $\Bin^\omega$ is !!{nonenumerable}
(\olref[nen]{thm:nonenum-bin-omega}). Daarom is $\Pow{\PosInt}$
!!{nonenumerable}.
\end{proof}

\begin{explain}
De richting van een reductie haal je makkelijk door elkaar. Neem
bijvoorbeeld een surjectieve functie $g \colon \Bin^\omega \to B$.
Daarmee bewijs je \emph{niet} dat $B$ !!{nonenumerable} is. (Neem $g
\colon \Bin^\omega \to \Bin$ met $g(s) = s(1)$, dus de functie die van
een rij met $0$'en en $1$'en het eerste !!{element} neemt. De functie
is !!{surjective}, want sommige rijen beginnen met $0$ en andere met
$1$. Maar $\Bin$ is eindig.) Ook de functie~$f$ moet !!{surjective}
zijn, anders werkt het argument niet. Zonder die eigenschap weten we
niet zeker of de lijst $f(x_1)$, $f(x_2)$, \dots{} alle
!!{element}en van~$B$ bevat.
Bijvoorbeeld:
\[
h(n) = \underbrace{000\dots0}_{\text{$n$ keer $0$}}
\]
Dit geeft een functie $h\colon \PosInt \to
\Bin^\omega$, maar $\PosInt$ is !!{enumerable}.
\end{explain}

\begin{prob}
Gebruik een reductie om te laten zien dat de verzameling van alle
\emph{verzamelingen van} paren positieve gehele getallen
!!{nonenumerable} is.
\end{prob}

\begin{prob}\label{sfr:siz:red:prob:nat-nat}
  Gebruik een reductie om te laten zien dat de verzameling~$X$ van alle
  functies $f\colon \Nat \to \Nat$ !!{nonenumerable} is (Aanwijzing:
  geef een surjectieve functie van $X$ naar~$\Bin^\omega$.)
\end{prob}

\begin{prob}
Laat met een reductie zien dat $\Nat^\omega$, de verzameling van alle
oneindige rijen natuurlijke getallen, !!{nonenumerable} is.
\end{prob}

\begin{prob}
Laat $P$ de verzameling zijn van functies van de positieve gehele
getallen naar de verzameling $\{0\}$. Laat $Q$ de verzameling zijn van
\emph{partiële} functies van de positieve gehele getallen naar de
verzameling $\{0\}$. Zulke partiële functies hoeven niet voor elke
invoer een waarde te hebben. Laat zien dat $P$~!!{enumerable} is en
$Q$~niet. (Aanwijzing: herleid het opsommen van $\Bin^\omega$ tot het
opsommen van~$Q$.)
\end{prob}

\begin{prob}
Laat $S$ de verzameling zijn van alle surjectieve functies van de
positieve gehele getallen naar de verzameling \{0,1\}. Met andere
woorden: $S$ bestaat uit alle surjectieve functies~$f\colon
\PosInt \to \Bin$. Laat zien dat $S$ !!{nonenumerable} is.
\end{prob}

\begin{prob}
Laat zien dat de verzameling~$\Real$ van alle reële getallen
!!{nonenumerable} is.
\end{prob}
\end{document}
